\section{Matroids and the self-dual reduction}\label{sec:matroids}

We first record the infinite-matroid conventions and explain why, for a
self-dual pair, it suffices to rule out an independent covering of the
whole ground set. This reduction is essential: a mixed packing/covering
partition need not be a packing or a covering on the whole set.

\subsection{Independence, closure, and duality}

An infinite matroid on a set $F$ is a family $\mathcal I\subseteq 2^F$
satisfying the following axioms
\cite[Section~1.1]{BDKPW2013}:
\begin{description}
\item[(I1)] $\varnothing\in\mathcal I$.
\item[(I2)] Every subset of a member of $\mathcal I$ belongs to $\mathcal I$.
\item[(I3)] If $I\in\mathcal I$ is not inclusion-maximal and
  $B\in\mathcal I$ is inclusion-maximal, there is an element
  $e\in B\setminus I$ such that $I\cup\{e\}\in\mathcal I$.
\item[(IM)] For every $I\in\mathcal I$ and $I\subseteq Y\subseteq F$,
  the family $\{J\in\mathcal I:I\subseteq J\subseteq Y\}$ has an
  inclusion-maximal member.
\end{description}
The inclusion-maximal independent sets are the \emph{bases}. In
particular, (IM) requires maximal extensions inside every subset, not
only inside $F$.

For a matroid $N=(F,\mathcal I)$ and $X\subseteq F$, its closure is
\[
 \operatorname{cl}_N(X)=X\cup
 \{e\in F:\text{there is an independent }I\subseteq X
                 \text{ with }I\cup\{e\}\text{ dependent}\}.
\]
This operator is extensive, monotone, and idempotent
\cite[Theorem~4.2]{BDKPW2013}. A set $S$ is \emph{spanning} if
$\operatorname{cl}_N(S)=F$.
Restrictions are matroids, and the complements of the bases of $N$
are the bases of its dual $N^*$
\cite[Theorems~3.1 and~3.4]{BDKPW2013}. We use
\[
 N.X=N/(F\setminus X)=(N^*\!\upharpoonright X)^*
\]
for contraction onto $X$. Self-duality will always mean the equality
$N=N^*$ on the same labeled ground set.

\begin{lemma}\label{lem:closure-duality}
Let $N$ be an infinite matroid on $F$.
\begin{enumerate}
\item If $S\subseteq X\subseteq F$, then
\[
 \operatorname{cl}_{N\upharpoonright X}(S)
   =\operatorname{cl}_N(S)\cap X.
\]
\item A subset $S\subseteq F$ is spanning if and only if it contains
  a basis of $N$.
\item A subset $S\subseteq F$ is spanning in $N$ if and only if
  $F\setminus S$ is independent in $N^*$.
\end{enumerate}
\end{lemma}
\begin{proof}
The first assertion follows directly from the closure definition:
for $e\in X$, independence of sets contained in $S\cup\{e\}$ is
the same in $N$ and in its restriction to $X$.

For the second, use (IM) to take a maximal independent subset $B$ of
$S$. For every $e\in S\setminus B$, the set $B\cup\{e\}$ is
dependent, so $S\subseteq\operatorname{cl}_N(B)$. Monotonicity and
idempotence therefore give
\[
 \operatorname{cl}_N(B)=\operatorname{cl}_N(S).
\]
If $S$ is spanning, so is $B$. An independent spanning set is a basis:
an independent extension by an element of its closure would contradict
the defining independent witness for closure. Conversely, a basis spans,
because adjoining any element outside it gives a dependent set; any set
containing a basis consequently spans.

If $S$ contains a basis $B$, then $F\setminus S$ is contained in
the dual basis $F\setminus B$. Conversely, extend an independent
$F\setminus S$ to a basis $B^*$ of $N^*$ using (IM).
The primal basis $F\setminus B^*$ is contained in $S$.
\end{proof}

\subsection{From a mixed partition to an independent covering}

\begin{lemma}\label{lem:selfdual-reduction}
Let $N_0,N_1$ be matroids on the same set $F$, with $N_i=N_i^*$
on the same labels. Suppose there is a partition
$F=P\mathbin{\dot\cup}C$, disjoint sets $S_0,S_1\subseteq P$
spanning $N_0\!\upharpoonright P,N_1\!\upharpoonright P$,
respectively, and independent sets $I_i$ of $N_i.C$ such that
$I_0\cup I_1=C$. Then $F$ can be covered by one independent set
of each of $N_0,N_1$.
\end{lemma}
\begin{proof}
Self-duality gives
\[
 N_i.C=(N_i\!\upharpoonright C)^*.
\]
Thus $T_i=C\setminus I_i$ spans $N_i\!\upharpoonright C$ by
Lemma~\ref{lem:closure-duality}. Also
$T_0\cap T_1=C\setminus(I_0\cup I_1)=\varnothing$.
Consequently the two sets
\[
 A_i=S_i\cup T_i\qquad(i=0,1)
\]
are disjoint. The restriction-closure identity shows that
$P\subseteq\operatorname{cl}_{N_i}(S_i)$ and
$C\subseteq\operatorname{cl}_{N_i}(T_i)$, so $A_i$ spans the
full matroid $N_i$. Its complement $F\setminus A_i$ is independent
in $N_i^*=N_i$. Finally,
\[
 (F\setminus A_0)\cup(F\setminus A_1)=F
\]
because $A_0\cap A_1=\varnothing$.
\end{proof}

We will construct a self-dual pair with no independent covering.
Lemma~\ref{lem:selfdual-reduction} will then exclude every mixed
packing/covering partition for that pair.
